\textbf{Phần 1. Trắc nghiệm lựa chọn}
\setcounterpageref{ex}{0}

\begin{ex} %[5H6N33]
Diện tích của một hình tam giác vuông có độ dài hai cạnh góc vuông lần lượt là $1{,}6\text{ dm}$ và $3{,}5\text{ dm}$ là
\choice
{$2{,}3\text{ dm}^2$}
{$5{,}6\text{ dm}^2$}
{$2{,}8\text{ m}^2$}
{\True $2{,}8\text{ dm}^2$}
\loigiai{
Diện tích hình tam giác vuông là: $\dfrac{1{,}6 \times 3{,}5}{2} = 2{,}8\text{ (dm}^2\text{)}$.\\
\textbf{Đáp án của câu hỏi là D}
}
\end{ex}

\begin{ex} %[5H6H33]
Diện tích của hình thang có độ dài hai đáy lần lượt là $14\text{ cm}$ và $36\text{ cm}$, chiều cao $17\text{ cm}$ là
\choice
{\True $425\text{ cm}^2$}
{$850\text{ cm}^2$}
{$300\text{ cm}^2$}
{$612\text{ cm}^2$}
\loigiai{
Diện tích hình thang là: $\dfrac{(14 + 36) \times 17}{2} = \dfrac{50 \times 17}{2} = 425\text{ (cm}^2\text{)}$.\\
\textbf{Đáp án của câu hỏi là A}
}
\end{ex}

\begin{ex} %[5H6H33]
Một tấm thảm hình tròn có bán kính $3\text{ dm}$. Chu vi và diện tích của tấm thảm đó lần lượt là
\choice
{$9{,}42\text{ dm}$ và $28{,}26\text{ dm}^2$}
{$18{,}84\text{ dm}$ và $18{,}84\text{ dm}^2$}
{\True $18{,}84\text{ dm}$ và $28{,}26\text{ dm}^2$}
{$28{,}26\text{ dm}$ và $28{,}26\text{ dm}^2$}
\loigiai{
Chu vi tấm thảm là: $3{,}14 \times 2 \times 3 = 18{,}84\text{ (dm)}$.\\
Diện tích tấm thảm là: $3{,}14 \times 3 \times 3 = 28{,}26\text{ (dm}^2\text{)}$.\\
\textbf{Đáp án của câu hỏi là C}
}
\end{ex}

\begin{ex} %[5H6H33]
Bánh xe lăn tay dành cho người khuyết tật có đường kính $50\text{ cm}$. Nếu bánh xe lăn trên mặt đất được $1\,000$ vòng thì người đi xe đó đã đi được quãng đường là
\choice
{$157\text{ m}$}
{\True $1\,570\text{ m}$}
{$314\text{ m}$}
{$3\,140\text{ m}$}
\loigiai{
Chu vi bánh xe là: $3{,}14 \times 50 = 157\text{ (cm)} = 1{,}57\text{ (m)}$.\\
Quãng đường đi được là: $1{,}57 \times 1\,000 = 1\,570\text{ (m)}$.\\
\textbf{Đáp án của câu hỏi là B}
}
\end{ex}

\begin{ex} %[5H6H33]
Một mặt bàn hình tròn có đường kính $1{,}4\text{ m}$. Diện tích của mặt bàn đó là
\choice
{$6{,}1544\text{ m}^2$}
{$4{,}396\text{ m}^2$}
{\True $1{,}5386\text{ m}^2$}
{$2{,}198\text{ m}^2$}
\loigiai{
Bán kính mặt bàn là: $1{,}4 : 2 = 0{,}7\text{ (m)}$.\\
Diện tích mặt bàn là: $3{,}14 \times 0{,}7 \times 0{,}7 = 1{,}5386\text{ (m}^2\text{)}$.\\
\textbf{Đáp án của câu hỏi là C}
}
\end{ex}

\begin{ex} %[5H6H33]
Một mảnh tôn hình tam giác có diện tích $72\text{ dm}^2$ và chiều cao tương ứng là $9\text{ dm}$. Độ dài cạnh đáy của mảnh tôn đó là
\choice
{$8\text{ dm}$}
{$32\text{ dm}$}
{$4\text{ dm}$}
{\True $16\text{ dm}$}
\loigiai{
Độ dài cạnh đáy là: $72 \times 2 : 9 = 16\text{ (dm)}$.\\
\textbf{Đáp án của câu hỏi là D}
}
\end{ex}

\begin{ex} %[5H6V33]
Mảnh đất của bác Tư dạng hình thang vuông có chiều cao $12\text{ m}$, đáy bé $18\text{ m}$ và đáy lớn bằng $\dfrac{4}{3}$ đáy bé. Bác Tư đã dành phần đất hình tam giác để hiến mở rộng đường (cắt theo chiều cao hạ từ đỉnh đáy bé xuống đáy lớn). Diện tích phần đất bác Tư đã hiến là
\choice
{$72\text{ m}^2$}
{\True $36\text{ m}^2$}
{$216\text{ m}^2$}
{$252\text{ m}^2$}
\loigiai{
Đáy lớn của mảnh đất là: $18 \times \dfrac{4}{3} = 24\text{ (m)}$.\\
Độ dài đáy của phần đất hình tam giác bị cắt đi là: $24 - 18 = 6\text{ (m)}$.\\
Diện tích phần đất bác Tư đã hiến là: $\dfrac{6 \times 12}{2} = 36\text{ (m}^2\text{)}$.\\
\textbf{Đáp án của câu hỏi là B}
}
\end{ex}

\begin{ex} %[5H6V33]
Cũng trong bài toán trên, diện tích phần đất còn lại của bác Tư sau khi đã hiến mở đường là
\choice
{\True $216\text{ m}^2$}
{$144\text{ m}^2$}
{$252\text{ m}^2$}
{$288\text{ m}^2$}
\loigiai{
Phần đất còn lại là hình chữ nhật có chiều dài $18\text{ m}$ và chiều rộng $12\text{ m}$.\\
Diện tích phần đất còn lại là: $18 \times 12 = 216\text{ (m}^2\text{)}$.\\
\textbf{Đáp án của câu hỏi là A}
}
\end{ex}

\begin{ex} %[5H6V33]
Một biển báo cấm đi ngược chiều là hình tròn có bán kính $35\text{ cm}$, phần hình chữ nhật màu trắng ở giữa có kích thước $50\text{ cm} \times 12\text{ cm}$. Diện tích phần màu đỏ của tấm biển đó là
\choice
{$3\,846{,}5\text{ cm}^2$}
{\True $3\,246{,}5\text{ cm}^2$}
{$600\text{ cm}^2$}
{$3\,446{,}5\text{ cm}^2$}
\loigiai{
Diện tích cả hình tròn biển báo: $3{,}14 \times 35 \times 35 = 3\,846{,}5\text{ (cm}^2\text{)}$.\\
Diện tích phần chữ nhật màu trắng: $50 \times 12 = 600\text{ (cm}^2\text{)}$.\\
Diện tích phần màu đỏ: $3\,846{,}5 - 600 = 3\,246{,}5\text{ (cm}^2\text{)}$.\\
\textbf{Đáp án của câu hỏi là B}
}
\end{ex}

\begin{ex} %[5H6V33]
Một thửa ruộng hình thang có diện tích $360\text{ m}^2$, chiều cao $15\text{ m}$. Biết đáy lớn hơn đáy bé $8\text{ m}$. Độ dài của đáy lớn thửa ruộng đó là
\choice
{$20\text{ m}$}
{$24\text{ m}$}
{\True $28\text{ m}$}
{$48\text{ m}$}
\loigiai{
Tổng độ dài hai đáy là: $360 \times 2 : 15 = 48\text{ (m)}$.\\
Độ dài đáy lớn là: $(48 + 8) : 2 = 28\text{ (m)}$.\\
\textbf{Đáp án của câu hỏi là C}
}
\end{ex}

\begin{ex} %[5H6V33]
Một bồn hoa gồm một hình vuông cạnh $4\text{ m}$ và $4$ nửa hình tròn gắn vào $4$ cạnh của hình vuông đó. Chu vi của bồn hoa đó là
\choice
{\True $25{,}12\text{ m}$}
{$12{,}56\text{ m}$}
{$50{,}24\text{ m}$}
{$16\text{ m}$}
\loigiai{
Chu vi bồn hoa gồm độ dài của $4$ nửa đường tròn đường kính $4\text{ m}$ (tương đương $2$ đường tròn đầy đủ).\\
Chu vi bồn hoa là: $2 \times (3{,}14 \times 4) = 25{,}12\text{ (m)}$.\\
\textbf{Đáp án của câu hỏi là A}
}
\end{ex}

\begin{ex} %[5H6V33]
Một hồ nước hình tròn có chu vi $188{,}4\text{ m}$. Xung quanh hồ người ta làm một con đường chạy viền rộng $2\text{ m}$. Diện tích của con đường chạy viền đó là
\choice
{$389{,}36\text{ m}^2$}
{$376{,}8\text{ m}^2$}
{$401{,}92\text{ m}^2$}
{\True $389{,}36\text{ m}^2$}
\loigiai{
Bán kính hồ nước là: $188{,}4 : 3{,}14 : 2 = 30\text{ (m)}$.\\
Bán kính của cả hồ và con đường là: $30 + 2 = 32\text{ (m)}$.\\
Diện tích con đường là: $3{,}14 \times (32 \times 32 - 30 \times 30) = 3{,}14 \times (1\,024 - 900) = 3{,}14 \times 124 = 389{,}36\text{ (m}^2\text{)}$.\\
\textbf{Đáp án của câu hỏi là D}
}
\end{ex}

\textbf{Phần 2. Trắc nghiệm đúng sai}
\setcounterpageref{ex}{0}

\begin{ex} %[5H6H33]
Xét tính đúng sai của các khẳng định sau:
\choiceTF
{\True Diện tích hình thang có hai đáy $14\text{ cm}, 36\text{ cm}$ và chiều cao $17\text{ cm}$ là $425\text{ cm}^2$.}
{Chu vi của tấm thảm hình tròn bán kính $3\text{ dm}$ là $9{,}42\text{ dm}$.}
{\True Diện tích tấm thảm hình tròn bán kính $3\text{ dm}$ là $28{,}26\text{ dm}^2$.}
{\True Bánh xe đường kính $50\text{ cm}$ lăn $1\,000$ vòng đi được quãng đường $1\,570\text{ m}$.}
\loigiai{
\begin{itemize}
    \item $S = \dfrac{(14 + 36) \times 17}{2} = 425\text{ cm}^2$.
    \item Chu vi: $3{,}14 \times 2 \times 3 = 18{,}84\text{ dm} \neq 9{,}42\text{ dm}$.
    \item Diện tích: $3{,}14 \times 3 \times 3 = 28{,}26\text{ dm}^2$.
    \item Quãng đường: $3{,}14 \times 0{,}5 \times 1\,000 = 1\,570\text{ m}$.
\end{itemize}
\begin{itemize}
    \item a) \textbf{Đúng}.
    \item b) \textbf{Sai}.
    \item c) \textbf{Đúng}.
    \item d) \textbf{Đúng}.
\end{itemize}
}
\end{ex}

\begin{ex} %[5H6V33]
Một tấm biển báo hình tròn có bán kính $35\text{ cm}$, trên mặt biển có dán một băng giấy hình chữ nhật màu trắng kích thước $50\text{ cm} \times 12\text{ cm}$, phần còn lại sơn màu đỏ. Xét tính đúng sai của các khẳng định sau:
\choiceTF
{\True Diện tích toàn bộ tấm biển hình tròn là $3\,846{,}5\text{ cm}^2$.}
{\True Diện tích của băng giấy hình chữ nhật màu trắng là $600\text{ cm}^2$.}
{Diện tích phần sơn màu đỏ là $3\,446{,}5\text{ cm}^2$.}
{\True Diện tích phần màu đỏ gấp hơn $5$ lần diện tích phần màu trắng.}
\loigiai{
\begin{itemize}
    \item $S_{tròn} = 3{,}14 \times 35 \times 35 = 3\,846{,}5\text{ cm}^2$.
    \item $S_{HCN} = 50 \times 12 = 600\text{ cm}^2$.
    \item $S_{đỏ} = 3\,846{,}5 - 600 = 3\,246{,}5\text{ cm}^2 \neq 3\,446{,}5\text{ cm}^2$.
    \item Tỉ số: $3\,246{,}5 : 600 \approx 5{,}41 > 5$.
\end{itemize}
\begin{itemize}
    \item a) \textbf{Đúng}.
    \item b) \textbf{Đúng}.
    \item c) \textbf{Sai}: Diện tích phần đỏ là $3\,246{,}5\text{ cm}^2$.
    \item d) \textbf{Đúng}.
\end{itemize}
}
\end{ex}

\textbf{Phần 3. Trả lời ngắn}
\setcounterpageref{ex}{0}

\begin{ex} %[5H6H33]
Tam giác vuông có hai cạnh góc vuông dài $1{,}6\text{ dm}$ và $3{,}5\text{ dm}$. Tính diện tích tam giác theo đề-xi-mét vuông. (Nhập kết quả dưới dạng số thập phân).
\shortans{2,8}
\loigiai{
\textbf{Đáp số:} $2{,}8$\\
$\dfrac{1{,}6 \times 3{,}5}{2} = 2{,}8\text{ (dm}^2\text{)}$.
}
\end{ex}

\begin{ex} %[5H6H33]
Hình thang có hai đáy dài $14\text{ cm}$ và $36\text{ cm}$, chiều cao $17\text{ cm}$. Tính diện tích hình thang theo xăng-ti-mét vuông.
\shortans{425}
\loigiai{
\textbf{Đáp số:} $425$\\
$\dfrac{(14 + 36) \times 17}{2} = 425\text{ (cm}^2\text{)}$.
}
\end{ex}

\begin{ex} %[5H6H33]
Mảnh tôn tam giác có diện tích $72\text{ dm}^2$ và chiều cao $9\text{ dm}$. Tính độ dài cạnh đáy theo đề-xi-mét.
\shortans{16}
\loigiai{
\textbf{Đáp số:} $16$\\
$72 \times 2 : 9 = 16\text{ (dm)}$.
}
\end{ex}

\begin{ex} %[5H6V33]
Bác Tư có mảnh đất hình thang vuông cao $12\text{ m}$, đáy bé $18\text{ m}$, đáy lớn $24\text{ m}$. Bác cắt bớt tam giác có đáy $6\text{ m}$ và chiều cao $12\text{ m}$ để hiến đất. Tính diện tích phần đất còn lại theo mét vuông.
\shortans{216}
\loigiai{
\textbf{Đáp số:} $216$\\
$18 \times 12 = 216\text{ (m}^2\text{)}$.
}
\end{ex}

\begin{ex} %[5H6V33]
Biển báo hình tròn bán kính $35\text{ cm}$ chứa hình chữ nhật $50\text{ cm} \times 12\text{ cm}$. Tính diện tích phần còn lại theo xăng-ti-mét vuông. (Nhập kết quả dưới dạng số thập phân).
\shortans{3246,5}
\loigiai{
\textbf{Đáp số:} $3246{,}5$\\
$3{,}14 \times 35 \times 35 - 50 \times 12 = 3\,846{,}5 - 600 = 3\,246{,}5\text{ (cm}^2\text{)}$.
}
\end{ex}

\begin{ex} %[5H6V33]
Mặt bàn hình tròn có đường kính $1{,}4\text{ m}$. Tính diện tích mặt bàn theo mét vuông. (Nhập kết quả dưới dạng số thập phân).
\shortans{1,5386}
\loigiai{
\textbf{Đáp số:} $1{,}5386$\\
$3{,}14 \times 0{,}7 \times 0{,}7 = 1{,}5386\text{ (m}^2\text{)}$.
}
\end{ex}
